\documentclass[11pt]{article}

\usepackage[margin=1in]{geometry}
\usepackage{amsmath,amssymb,amsthm}
\usepackage{graphicx}
\usepackage[hidelinks]{hyperref}

\newtheorem{theorem}{Theorem}
\newtheorem{lemma}{Lemma}
\newtheorem{remark}{Remark}

\newcommand{\C}{\mathbb C}
\newcommand{\N}{\mathbb N}
\newcommand{\HH}{\mathcal H}
\newcommand{\BB}{\mathcal B}
\newcommand{\norm}[1]{\left\lVert #1 \right\rVert}
\newcommand{\bprop}{\textup{(}\beta\textup{)}}

\title{A Scalar-Shift Counterexample to the Local Bishop-\(\beta\) Sum Question}
\author{\texttt{agent\_lane\_03}}
\date{2026-07-18}

\begin{document}
\maketitle

\section*{Source and Question}

The source paper is Salah Mecheri, \emph{On the stability of the Bishop's
property(\(\beta\)) under compact perturbations}, arXiv:2510.09289.  On page
5, immediately after Theorem 2.1, it asks the following localized question:

\begin{quote}
If \(S,T\in B(X)\) have Bishop's property \((\beta)\) at \(\lambda_0\), does
\(S+T\) have Bishop's property \((\beta)\) at \(\lambda_0\)?
\end{quote}

\begin{figure}[h]
\centering
\includegraphics[width=0.96\linewidth]{figures/open_problem_crop.png}
\caption{The localized sum question in arXiv:2510.09289, p. 5.}
\end{figure}

The answer to this question, as stated with the same fixed base point
\(\lambda_0\), is negative.

\section*{Definitions}

Let \(A\in\BB(X)\), where \(X\) is a complex Banach space.  We say that \(A\)
has Bishop's property \(\bprop\) at \(\lambda_0\) if there is \(r>0\) such
that, for every open set \(U\subset D(\lambda_0,r)\) and every sequence
\((f_n)\subset\mathcal O(U;X)\),
\[
   (A-\lambda I)f_n(\lambda)\longrightarrow 0
   \quad\hbox{in }\mathcal O(U;X)
\]
implies
\[
   f_n(\lambda)\longrightarrow 0
   \quad\hbox{in }\mathcal O(U;X).
\]
The single-valued extension property, abbreviated SVEP, at \(\lambda_0\) says
that the only analytic \(f\) near \(\lambda_0\) with
\((A-\lambda I)f(\lambda)=0\) is \(f\equiv 0\).

\section*{Proof Intuition}

The localized formulation is not stable under scalar translation at a fixed
base point.  Start from an operator \(B\) that fails Bishop's property
\(\bprop\) at \(0\).  Shift it far away to \(B-2I\); now \(0\) lies in the
resolvent set, so the shifted operator has \(\bprop\) at \(0\).  The scalar
operator \(2I\) also has \(\bprop\) at \(0\).  But adding the two commuting
operators gives back \(B\), which was chosen to fail \(\bprop\) at \(0\).

The backward unilateral shift on \(\ell^2\) supplies the required \(B\): it has
a nonzero analytic eigenvector field \(z\mapsto(1,z,z^2,\ldots)\) for
\(|z|<1\), so it fails SVEP, and hence fails Bishop's property \(\bprop\), at
the origin.

\section*{The Counterexample}

\begin{lemma}\label{lem:resolvent}
If \(\lambda_0\in\rho(A)\), then \(A\) has Bishop's property \(\bprop\) at
\(\lambda_0\).
\end{lemma}

\begin{proof}
Choose \(r>0\) such that \(\overline{D(\lambda_0,r)}\subset\rho(A)\).  Let
\(U\subset D(\lambda_0,r)\) be open and suppose
\((A-\lambda I)f_n(\lambda)\to0\) in \(\mathcal O(U;X)\).  For each compact
\(K\subset U\), the resolvent is uniformly bounded on \(K\), so
\[
  \sup_{\lambda\in K}\norm{f_n(\lambda)}
  \le
  \sup_{\lambda\in K}\norm{(A-\lambda I)^{-1}}\,
  \sup_{\lambda\in K}\norm{(A-\lambda I)f_n(\lambda)}
  \longrightarrow 0.
\]
Thus \(f_n\to0\) in \(\mathcal O(U;X)\).
\end{proof}

\begin{lemma}\label{lem:beta-svep}
If \(A\) has Bishop's property \(\bprop\) at \(\lambda_0\), then \(A\) has
SVEP at \(\lambda_0\).
\end{lemma}

\begin{proof}
If \(f\) is a nonzero analytic solution of
\((A-\lambda I)f(\lambda)=0\) on a neighborhood of \(\lambda_0\), then the
constant sequence \(f_n=f\) satisfies
\((A-\lambda I)f_n(\lambda)\equiv0\), but \(f_n\) does not converge to \(0\).
This contradicts Bishop's property \(\bprop\).
\end{proof}

\begin{theorem}
There are commuting operators \(S,T\in\BB(\ell^2(\N_0))\) such that both \(S\)
and \(T\) have Bishop's property \(\bprop\) at \(0\), but \(S+T\) does not have
Bishop's property \(\bprop\) at \(0\).
\end{theorem}

\begin{proof}
Let \(B\) be the backward unilateral shift on
\(\HH=\ell^2(\N_0)\):
\[
   B(x_0,x_1,x_2,\ldots)=(x_1,x_2,x_3,\ldots).
\]
Define
\[
   S=2I,\qquad T=B-2I.
\]
Plainly \(S\) and \(T\) commute.

The operator \(S\) has Bishop's property \(\bprop\) at \(0\), since \(0\in
\rho(S)\).  The operator \(T\) also has Bishop's property \(\bprop\) at \(0\).
Indeed, for \(|\lambda|<1/2\),
\[
  T-\lambda I = B-(2+\lambda)I
\]
is invertible by the Neumann series, because
\[
   \left\|\frac{B}{2+\lambda}\right\|
   \le \frac{1}{|2+\lambda|}
   < 1.
\]
Hence \(0\in\rho(T)\), and Lemma~\ref{lem:resolvent} applies.

It remains to show that \(S+T=B\) does not have Bishop's property \(\bprop\)
at \(0\).  For \(|z|<1\), put
\[
   g(z)=(1,z,z^2,\ldots)\in\ell^2(\N_0).
\]
The map \(g:D(0,1)\to\ell^2(\N_0)\) is analytic and nonzero, and
\[
   Bg(z)=(z,z^2,z^3,\ldots)=z\,g(z).
\]
Thus
\[
   (B-zI)g(z)=0\qquad(|z|<1),
\]
so \(B\) fails SVEP at \(0\).  By Lemma~\ref{lem:beta-svep}, \(B\) cannot
have Bishop's property \(\bprop\) at \(0\).  Since \(S+T=B\), this proves the
claim.
\end{proof}

\section*{Scope and Limitations}

This is a counterexample to the localized question exactly as printed on page
5 of the source paper: both summands have Bishop's property \(\bprop\) at the
same fixed point \(0\), while their sum fails it at that same point.

It does not settle the broader global problem discussed in the source
introduction, namely whether sums or products of commuting operators that have
Bishop's property \(\bprop\) everywhere must again have Bishop's property
\(\bprop\) everywhere.  In the example above, \(T=B-2I\) is used only through
its local behavior at \(0\); no global Bishop-\(\beta\) assertion is made for
\(T\).

\section*{Novelty Check}

The run's cheap indexes were searched for arXiv:2510.09289, the source title,
``Bishop's property'', ``Bishop beta'', ``localized Bishop'', ``SVEP'',
``commuting perturbation'', and ``backward shift''.  No prior packet, proof
gap, or attempt for this exact localized question was found.

A bounded web search for the exact source-question phrase and nearby phrases
found the source preprint, the older/global formulation in Mecheri's 2012
paper, and standard background statements such as Bishop's property
\(\bprop\) implying SVEP.  It did not reveal an existing scalar-shift
counterexample to the fixed-base-point localized sentence.  This is not an
exhaustive literature search; because the example is elementary, it may be
folklore.  The mathematical claim above is nevertheless self-contained.

\section*{Verification Notes and Human Review}

Recommended review verdict: likely valid counterexample to the localized
question.  The key checks are:

\begin{itemize}
\item verify that the source question asks for the sum at the same fixed
      \(\lambda_0\);
\item verify Lemma~\ref{lem:resolvent}, especially the use of the resolvent
      neighborhood for a local property;
\item verify that \(z\mapsto(1,z,z^2,\ldots)\) is a nonzero analytic
      \(\ell^2\)-valued solution of \((B-zI)g(z)=0\) near \(0\).
\end{itemize}

No computational verification is needed.

\begin{thebibliography}{9}

\bibitem{Mecheri2025}
Salah Mecheri,
\emph{On the stability of the Bishop's property(\(\beta\)) under compact
perturbations},
arXiv:2510.09289, 2025.

\bibitem{LaursenNeumann}
K. B. Laursen and M. M. Neumann,
\emph{An Introduction to Local Spectral Theory},
London Mathematical Society Monographs, Clarendon Press, Oxford, 2000.

\end{thebibliography}

\end{document}
